\PassOptionsToPackage{dvipsnames,svgnames,table}{xcolor}
\documentclass[11pt]{article}
\usepackage[a4paper,margin=22mm]{geometry}
\usepackage[no-math]{fontspec}
\setmainfont{DejaVu Serif}
\usepackage{amsmath,amssymb,amsthm,mathtools,graphicx,xcolor,hyperref,xurl,xspace,ifthen,enumerate}
\usepackage[shortlabels]{enumitem}
\setlength{\emergencystretch}{4em}
\AtBeginDocument{\special{pdf:mapline pzdr Dingbats <uzdr.pfb}}
\SetMathAlphabet{\mathbf}{normal}{OT1}{cmr}{bx}{n}
\SetMathAlphabet{\mathbf}{bold}{OT1}{cmr}{bx}{n}
\newtheorem{theo}{Theorem}[section]
\newtheorem{lemm}[theo]{Lemma}
\newtheorem{thm}[theo]{Theorem}
\newtheorem{proposition}[theo]{Proposition}
\newtheorem{corollary}[theo]{Corollary}
\newtheorem{definition}[theo]{Definition}
\newtheorem{theorem}[theo]{Theorem}
\newtheorem*{remas*}{Remarks}
\newtheorem{remas}[theo]{Remarks}
\newtheorem*{exam*}{Example}
\newtheorem{ques}[theo]{Question}
\providecommand{\og}{«\nobreakspace}
\providecommand{\fg}{\nobreakspace»}
\newtheorem{lem}[theo]{Lemma}
\newtheorem{lemma}[theo]{Lemma}
\newtheorem{prop}[theo]{Proposition}
\newtheorem{coro}[theo]{Corollary}
\newtheorem{corol}[theo]{Corollary}
\newtheorem*{conj*}{Conjecture}
\newtheorem{conj}[theo]{Conjecture}
\newtheorem{conjecture}[theo]{Conjecture}
\newtheorem{Proposition}[theo]{Proposition}
\newtheorem{theoreme}[theo]{Theorem}
\theoremstyle{definition}
\newtheorem{defi}[theo]{Definition}
\newtheorem{exam}[theo]{Example}
\newtheorem{exem}[theo]{Example}
\newtheorem{remark}[theo]{Remark}
\newtheorem{example}[theo]{Example}
\newtheorem{exams}[theo]{Examples}
\newtheorem{rema}[theo]{Remark}
\newtheorem{nota}[theo]{Notation}
\newtheorem{notation}[theo]{Notation}
\newtheorem*{notas}{Notations}
\newtheorem{result}[theo]{Result}
\newtheorem{question}[theo]{Question}
\newtheorem{prob}[theo]{Problem}
\newtheorem*{rema*}{Remark}
\newtheorem*{defi*}{Definition}
\newtheorem*{theo*}{Theorem}
\newtheorem*{lemm*}{Lemma}
\newtheorem*{prop*}{Proposition}
\newcommand{\enoncename}{}
\newtheorem{corpusenonce}[theo]{\enoncename}
\newenvironment{enonce}[2][]{\renewcommand{\enoncename}{#2}\begin{corpusenonce}}{\end{corpusenonce}}
\newenvironment{enonce*}[2][]{\par\medskip\noindent\textbf{#2.}\itshape}{\par\medskip}
\newenvironment{proofc}[1][\proofname]{\begin{proof}[#1]}{\end{proof}}
\newcommand{\Mc}{\mathcal{M}}
\newcommand{\pointir}{.\ ---\ }
\newcommand{\equalenv}[2]{\expandafter\let\csname#1\expandafter\endcsname\csname#2\endcsname\expandafter\let\csname end#1\expandafter\endcsname\csname end#2\endcsname}
\providecommand{\diagbox}[3][]{\begin{tikzpicture}[x=1em,y=1em,baseline=(current bounding box.center)]\useasboundingbox(0,0) rectangle(3,1.4);\draw(0,1.4)--(3,0);\node at(.5,.3){#2};\node at(2.5,1.1){#3};\end{tikzpicture}}
\providecommand{\eg}{e.g.\ }
\providecommand{\ie}{i.e.\ }
\providecommand{\cf}{cf.\ }
\providecommand{\longthanks}[1]{\par\smallskip\noindent #1\par\smallskip}
\providecommand{\qbinom}[2]{\genfrac{[}{]}{0pt}{}{#1}{#2}}

\providecommand{\mathof}[1]{\mathrm{#1}}

\numberwithin{equation}{section}


\def\endappendix{}


\usepackage{latexsym,amsmath,amssymb,verbatim,tikz,tikz-cd}
\usepackage{mathtools}
\usepackage{graphicx,xcolor}

\usepackage[english]{babel}
\usepackage{amssymb}
\usepackage{amsmath}
\usepackage{amsthm}
\usepackage{graphicx}
\usepackage{csquotes}
\usepackage{bm}
\usepackage{enumitem}
\usepackage{subcaption}
\usepackage{caption}
\renewcommand{\thesubfigure}{\Alph{subfigure}}
\usepackage{dirtytalk}
\providecommand{\nopunct}{\spacefactor1007\relax}
\newtheorem{problem}[theo]{Problem}

\newcommand{\sym}{\mathrm{Sym}}
\newcommand{\qsym}{\mathrm{QSym}}
\newcommand{\Des}{\operatorname{Des}}
\newcommand{\Asc}{\operatorname{Asc}}
\newcommand{\Syt}{\operatorname{SYT}}
\newcommand{\Short}{\operatorname{Short}}
\newcommand{\Loop}{\operatorname{Loop}}
\newcommand{\Stable}{\operatorname{Stable}}
\newcommand{\len}{\operatorname{len}}
\newcommand{\height}{\operatorname{height}}
\newcommand{\core}{\operatorname{core}}
\newcommand{\res}{\operatorname{res}}
\newcommand{\row}{\operatorname{row}}
\newcommand{\insrt}{\operatorname{insert}}
\newcommand{\UP}{\operatorname{UP}}
\newcommand{\DOWN}{\operatorname{DOWN}}


\newcommand{\ZZ}{\mathbb{Z}}
\newcommand{\NN}{\mathbb{N}}
\newcommand{\RR}{\mathbb{R}}
\newcommand{\QQ}{\mathbb{Q}}
\newcommand{\PP}{\mathbb{P}}

\newcommand{\F}{\mathcal{F}}
\newcommand{\Q}{\mathcal{Q}}
\newcommand{\A}{\mathcal{A}}
\newcommand{\M}{\mathcal{M}}
\newcommand{\I}{\mathcal{I}}
\newcommand{\PM}{\mathcal{PM}}
\newcommand{\Pcal}{\mathcal{P}}


\newcommand{\odds}{\mathrm{odds}}

\makeatletter
\newcommand{\proofstep}[1]{
  \par
  \addvspace{\medskipamount}
  \textbf{#1\@addpunct{.}}\enspace\ignorespaces
}

\newcommand\restr[2]{{
  \left.\kern-\nulldelimiterspace 
  #1 
  \vphantom{\big|} 
  \right|_{#2} 
  }}
\makeatother












\title{Three structural properties of the stability relation: intersecting chords, unstable chords, and invariant blocks}
\author{Avichai Marmor}
\date{doi:10.5802/alco.351; CC BY 4.0}
\begin{document}
\maketitle
\noindent\textit{Editorial note.} The counted statement is the author's labelled lemma `lem:properties of core` (source lines 873-881); the proof environment begins on the next line after \textbackslash{}end\{lemma\} apart from one blank line, so the pairing is by adjacency.\par\medskip
\noindent\textit{Editorial note.} The retained context is the author's definition of an m-invariant set, which the third part of the statement and its proof invoke.\par\medskip
\noindent\textit{Editorial note.} The retained preamble is the editorially prepared standalone preamble retained for this source.\par\medskip
\noindent\textit{Editorial note.} Every label used inside the retained text is defined inside the retained spans (the self-reference in the enumeration options is to the lemma's own label), so the package declares no external reference.\par\medskip
\noindent\textit{Editorial scope.} One original result of the pinned publisher article, counted once: the structural properties of the stability relation for matchings (author label lem:properties of core).\par\medskip
\noindent\textit{Allowed standard prerequisites (editorial).}\begin{itemize}
\item Stable and unstable chords. A chord is stable if it belongs to the core, equivalently if its endpoints are not removed during the reduction process; otherwise it is unstable.. The three parts are statements about stability and are proved by tracking which chords can be removed.
\item Invariant set and restricted matching. A set S is m-invariant if every chord of m meets S in either two or zero vertices, and res\_S(m) is the restriction of m to S.. The third part assumes [i,j] is m-invariant and restricts the matching to it.
\item Short chords and reduction. The reduction process repeatedly removes short chords, i.e. chords of the form (i,i+1), together with their endpoints.. Both the first and second parts argue about which chords can become short.
\end{itemize}\medskip
\begin{definition}
\label{def:restriction of matching}
    Given a matching $m \in \M_N$, we say that a set $S \subseteq [N]$ is $m$-invariant if for every chord $(i,j) \in m$ we have $i \in S$ if and only if $j \in S$. Moreover, for a matching $m$ and an $m$-invariant set $S$ we define the \emph{restriction} of $m$ to $S$, denoted $\res_S(m)$, to be a matching on $S$ which is obtained by removing all the vertices not in $S$ from $m$.
\end{definition}

\begin{lemma}
\label{lem:properties of core}
    The reduction process has the following properties:
    \begin{enumerate}[label=(\arabic*), ref=\ref{lem:properties of core}\text{ (part~}\arabic*)]
        \item \label{lem:properties of core--intersect stable} If a chord intersects another chord then it is stable.
        \item \label{lem:properties of core--chord unstable iff all midpoints unstable} A chord $(i,j)$ is stable if and only if there exists a stable vertex in $\{i+1,\dots,j-1\}$.
        \item \label{lem:properties of core--perfect non-crossing unstable} Given a matching $m$ and $i < j$, if the set $[i,j]$ is $m$-invariant \textup{(}as defined in Definition~\ref{def:restriction of matching}\textup{)} and the restricted matching $\res_{[i,j]}(m)$ is perfect and non-crossing, then $\ell \notin \Stable(m)$ for all $\ell \in [i,j]$.
    \end{enumerate}
\end{lemma}

\begin{proof} \hfill %
    \begin{enumerate}
        \item Assume that $(i_1,i_3),\; (i_2,i_4) \in m$ for $i_1 < i_2 < i_3 < i_4$. As long as the chord $(i_2,i_4)$ is not removed, the chord $(i_1,i_3)$ does not become short and cannot be removed. On the other hand, as long as the chord $(i_1,i_3)$ is not removed, the chord $(i_2,i_4)$ cannot be removed. Therefore, both chords cannot be removed during the reduction process.
        \item If the chord $(i,j)$ is unstable, then after removing some unstable vertices it becomes short, implying that every vertex between $i$ and $j$ is unstable. On the other hand, 
        if all the vertices between $i$ and $j$ are unstable, then they will eventually be removed, making the chord $(i,j)$ short, so the chord $(i,j)$ is unstable too.
        \item %
        Since the set $[i,j]$ is $m$-invariant and the restricted matching is perfect, Definition~\ref{def:restriction of matching} implies that for every $i \le i_1 \le j$ there exists $i \le i_2 \le j$ such that $i_1 \ne i_2$ and $(i_1,i_2) \in m$ or $(i_2,i_1) \in m$.
        Notice that the restricted matching $\res_S(m)$ where $S = \Stable(m) \cap [i,j]$ (\ie the matching that consists of the stable chords $(i_1,i_2) \in m$ with $i \le i_1<i_2 \le j$) is short-chord-free, non-crossing and perfect. The only such a matching is the empty matching $\emptyset \in \M_0$, so every $i \le i_1 \le j$ is unstable.
        \qedhere
    \end{enumerate}
\end{proof}

\end{document}
